\section{模型评价与推广}
\begin{enumerate}
\item \textbf{客观评价模型。} 从模型假设、计算效率、稳定性、适用范围等方面分析模型的优点与不足，做到客观真实。

\item \textbf{结合模型特点。} 所有评价都应结合本文建立的模型展开，避免使用“精度高”“效果好”等空泛描述。

\item \textbf{提出改进方向。} 对存在的问题可简要说明今后的优化思路，但不必重新建立新模型。

\item \textbf{分析推广价值。} 说明模型是否能够应用于类似问题，以及推广过程中需要调整的内容。
\end{enumerate}
\hspace*{\parindent}\textbf{本部分主要体现对模型的全面认识，不是论文重点，做到分析客观、内容精炼即可。}


\textbf{一、模型优点}

\begin{enumerate}
\item 分析模型是否具有较好的\textbf{科学性和合理性}，能够较真实地描述实际问题。

\item 说明模型是否具有\textbf{较好的可解释性}，变量和参数具有明确的实际意义。

\item 评价模型的\textbf{计算效率}，分析算法是否收敛较快、计算复杂度是否较低。

\item 说明模型是否具有\textbf{较好的稳定性}，当参数或数据发生小范围变化时，结果是否保持稳定。

\item 分析模型是否具有\textbf{较好的通用性}，是否能够适用于同类问题。
\end{enumerate}

\textbf{二、模型缺点}

\begin{enumerate}
\item 分析模型假设是否过于理想化，对实际问题进行了哪些简化。

\item 说明模型对数据质量、参数选取或初始条件是否较为敏感。

\item 分析模型适用范围是否有限，在哪些条件下可能不再适用。

\item 指出模型是否忽略了部分影响因素，或某些复杂关系未能充分考虑。

\item 结合实际情况说明模型仍存在的不足，并提出后续改进方向。
\end{enumerate}

\textbf{三、模型推广}

\begin{enumerate}
\item 说明模型的核心思想是否可以应用于同类型问题。

\item 分析模型推广到其他地区、行业或研究对象时，需要调整哪些参数或约束条件。

\item 若更换数据来源或增加新的影响因素，应说明模型需要进行哪些修改。

\item 可以结合人工智能、大数据、优化算法等方法，简要说明模型未来的改进和应用方向。
\end{enumerate}
\paragraph{写作提示}
\begin{enumerate}
\item 优点和缺点应一一对应，评价要有依据，不宜过分夸大或刻意贬低。

\item 推广应立足于模型本身，不宜脱离论文内容进行过多延伸。

\item 不必每一点都详细展开，可结合实际模型选择其中2～4点进行说明。

\item 建议控制篇幅，一般\textbf{0.5～1页}即可。
\end{enumerate}
\paragraph{写作原则：评价客观、优缺点对应、推广合理、内容简洁。}
\subsection{模型的优点}
\begin{enumerate}
\item 模型充分结合实际，简化xx，xx，xx条件，考虑了诸多重要因素得到合理的模型，如:xx，xx，xx。这样得到的模型贴合实际，具有较高的应用价值，可以推广到xx；【模型的假设好】
\item 模型运用xx和xx思想，抓住影响xx问题的重要因素，将复杂的xx问题转化为简单的xx问题，合理设置参数，模型的输出结果符合题目要求，能解决实际问题；【模型的参数好】
\item 本文使用的xx算法具有xx，xx，xx等优点，对于求解xx模型非常适用；【模型的求解算法好】
\item 本文得到的xx(安排方案、策略)具有效率高、输出稳定、xx均衡等特点，基本不存在xx，xx，xx等问题，在现有条件下能有效提高生产效率。【模型的结果好】
\end{enumerate}
\subsection{模型的不足}
\begin{enumerate}
\item 实际应用中，xx和xx可能也是重要的因素，但本文未能考虑到这些因素的影响，一定程度上影响了模型的准确性；【模型不够好】
\item 本文提出的模型对于现有条件使用效果较好，由于时间问题没有对其他情况进行检验。对于其他情形(如:xx，xx)，可能无法达到较好的效果；【模型适用范围较窄】
\item 实际上xx，xx的影响不一定是线性的，而本文将其作为线性因子处理，忽略了边际效应的影响。【模型的因素有些不好】
\end{enumerate}
\subsection{模型的推广}
【推广：在xx方面，可以将xx参数替换成xx参数，从而解决xx问题；
改进：结合参考文献xx，进一步考虑xx的影响，从而得到更合理的模型；
这部分不用太多，至少4行】
